\section{The four compatibility conditions after a Hecke cut}

% BEGIN T83C1
\paragraph{T83C1. Descent of the integral central action.}
Let $c$ be any cut of the original datum $D$. Thus $c$ specifies
the indices to keep and, at each kept $i$, a nonempty subset
$S\subseteq X_i$ of its finite set of Hecke points. Fix such an
$i$. Write $O$ for its coefficient DVR, $K$ for its fraction
field, $M$ for the integral Hecke module, $T\subset O^{X_i}$
for the separating Hecke algebra, and $\rho$ for its action.
The cut module and algebra of LN3 are
\[
 M'=M/M_{\rm out},\qquad
 M_{\rm out}=\bigcap_{f\in T,\ f|_{X_i\setminus S}=0}\ker\rho(f),
 \qquad T'=\operatorname{im}(T\longrightarrow O^S).
\]
Write $q$ for the quotient map and $q_K$ for its scalar extension.
An operator $b\in\operatorname{End}_O(M)$ commuting with $T$
preserves $M_{\rm out}$, because each $\rho(f)$ in this
intersection kills $bm$ when it kills $m$. Hence
$b'(q m)=q(bm)$ gives a unique integral endomorphism. For an
algebra of commuting-with-$T$ operators this assignment is an
$O$-algebra homomorphism: compute the two sides of each sum,
product, scalar and unit identity on the representatives $qm$.
In particular this applies to the given action $z_O$ of the
center of $O[G_{0g}]$ in a witness $P$ for the original
$\mathrm{Pend}$ record.

The action of $T$ itself on $M'$ factors through $T'$. If
$f|_S=g|_S$, then $h(f-g)=0$ for all $h$ supported on $S$,
so $(\rho(f)-\rho(g))m\in M_{\rm out}$. Thus $z'_O(z)$ commutes
with every element of $T'$, by lifting that element to $T$ and
testing on $qm$. Rationalizing gives, on pure tensors and hence
on all of $K\otimes_O M$,
\[
       q_K\rho_K(f)=\rho'_K(f|_S)q_K,\qquad
       q_K z_{O,K}(z)=z'_{O,K}(z)q_K .
\]
Every integral calculation was made on the given quotient, including
at primes dividing the order of the finite group.
% END T83C1

% BEGIN T83C2
\paragraph{T83C2. The projected types and their entire eigenspaces.}
For $x\in X_i$ let $V_x$ consist of the vectors in $K\otimes_O M$
on which each $f$ acts by $f(x)$. For $x\in S$ let $V'_x$ denote
the corresponding space in $K\otimes M'$. For each $x$ use
separation to form $h_x=\prod_{y\ne x}(f_{xy}-f_{xy}(y))$ as
in LN3. This has support $\{x\}$ and $h_x(x)\ne0$. If $x\notin S$,
then $h_x|_S=0$ and the first identity in T83C1 gives
$h_x(x)q_K V_x=0$. If $x\in S$, that same identity and
surjectivity $T\to T'$ give $q_K V_x\subseteq V'_x$.
For the reverse inclusion, take $u\in V'_x$ and choose any
$v\in K\otimes M$ with $q_Kv=u$. The vector
\[
               h_x(x)^{-1}\rho_K(h_x)v
\]
lies in $V_x$ because $fh_x=f(x)h_x$, and maps to $u$ since
$h_x|_S$ acts there by $h_x(x)$. The map on $V_x$ is injective
as well: for $x\in S$ the operator $\rho(h_x)$ kills
$M_{\rm out}$ and hence factors through $q$ as a map to $M$.
A vector of $V_x$ in $\ker q_K$ is therefore annihilated by
the nonzero scalar $h_x(x)$.

For a kept $i$, the cut reads the original point of an
automorphic representative $\Pi$ and retains its label exactly
when it belongs to $S$. Thus for $x\in S$ the set of representatives
over $x$ is exactly the original set over its underlying point.
The new projection of a type vector $a$ is
$q_K\operatorname{pr}_{M}(a)$. If the old point was absent its
old projection was zero. If the point was discarded, the old
projection belongs to the discarded eigenspace and is killed
by $q_K$. If its point is retained, the new projection lies in
$V'_x$. These three cases prove both clauses of the projected
type condition in the cut $\mathrm{Pend}$ record.

For the record's integral action comparison, suppose that the
complex type action of $z$ sends the vector corresponding to $a$
to that corresponding to $b$. The original comparison says
$\operatorname{pr}_M(b)=z_{O,K}(z)\operatorname{pr}_M(a)$.
Apply $q_K$ and T83C1 to obtain precisely the cut comparison.
The stability of the coefficient-field type space is the original
stability on that same space. The lists of local types, their
occurrence conditions, idempotent expansions and rationality,
the chosen primitive roots, the inertia characters and their two
characteristic-polynomial formulas carry over with the original
coefficient embeddings. The parameter-surjectivity clause carries
over too: a representative witnessing a base-change parameter at
a retained point is a representative with that parameter at its
underlying original point, so its original global parameter
witness has all the required local values. This constructs
$P'$ on the cut with the given local labels and integral
central action.
% END T83C2

% BEGIN T83C3
\paragraph{T83C3. Transport of the native four compatibility conditions.}
Suppose $\mathrm{Tie4}(D)$, with witness $P$ for its second
condition. For each retained $i$, each $x\in S$, each central
element $z\in Z(O[G_{0g}])$, and each $a\in O$, assume that
the complex action of $z$ is the coefficient image of $a$ on
every type vector of every representative over this cut point.
By T83C2 the condition concerns exactly the original representatives
and type spaces over the underlying $x$. The original
$\mathrm{EigSpan}(P)$ therefore says $z_{O,K}(z)v=av$ for
every $v\in V_x$. Given any $u\in V'_x$, lift it in $V_x$ by
T83C2 and apply T83C1. This gives
$z'_{O,K}(z)u=q_K(av)=au$, the full $\mathrm{EigSpan}(P')$.
The depth-zero condition $\mathrm{admDZ}$ tests the same local
types, their squarefree labels, the two excluded character roots,
monodromy and the local-parameter polynomial at each retained
index. Its original witness thus gives $\mathrm{admDZ}(P')$.

The other three conditions restrict at their index. In
$\mathrm{tieChi}$ use the same unit at each generator-prime
place and its prescribed inverse at the conjugate place.
In $\mathrm{TieInePoly}$ every continuous representation and
potentially crystalline period, polynomial and inertia element
tested at the retained index are in the original test with the
same coefficient embedding, so its polynomial action vanishes.
In $\mathrm{SerLoc}$ the series idempotent acts on the same
automorphic parameter space, giving the same polynomial label
from a nonzero action. Together these give
\[
           \mathrm{Tie4}(D)\quad\Longrightarrow\quad
                         \mathrm{Tie4}(\operatorname{cut}(D,c))
\]
for every $c$, before imposing any goodness condition on the cut.
This supplies the transported $\mathrm{Tie4}$ component of the
selected-family construction when the original four conditions
are given.
% END T83C3
